\section{Cronograma}
\label{sec:cronograma}

Como a data oficial do Seminário de Andamento ainda não foi registrada neste
repositório, o cronograma usa $S$ para representar a semana da apresentação. As
datas serão convertidas assim que houver comunicação institucional confirmada.

\begin{table}[H]
  \centering
  \caption{Cronograma relativo até o Seminário de Andamento}
  \label{tab:cronograma}
  \small
  \begin{tabularx}{\textwidth}{@{}L{2cm}Y@{}}
    \toprule
    \textbf{Período} & \textbf{Entregas verificáveis} \\
    \midrule
    $S-8$ a $S-7$ & Confirmar normas; completar metadados; congelar protocolo de revisão; preparar instrumentos, autorizações e solicitação de dados. \\
    $S-7$ a $S-5$ & Executar buscas; selecionar estudos; consolidar contextualização, problema, objetivos e fundamentação. \\
    $S-5$ a $S-3$ & Coletar dados autorizados; realizar entrevistas aprovadas; validar requisitos e arquitetura; consolidar os simuladores. \\
    $S-3$ a $S-2$ & Ampliar testes de componentes e integração; calcular primeira análise econômica; registrar limitações e resultados sem extrapolação. \\
    $S-2$ a $S-1$ & Submeter ao orientador; corrigir texto, referências, tabelas e compilação; fechar versão rastreável. \\
    $S-1$ a $S$ & Produzir apresentação modular, revisar o PDF, ensaiar e reservar margem para correções. \\
    \bottomrule
  \end{tabularx}
  \\[2pt]\fontepropria
\end{table}

Após o seminário, as atividades serão: finalizar bancada, concluir integração
autorizada se disponível, ampliar os testes, executar a avaliação com usuários,
fechar a análise econômica, revisar ameaças à validade e redigir implementação,
demonstração, resultados, conclusão e defesa final.

\pendente{inserir data oficial do seminário, duração, formato, limite de páginas
e marcos do orientador; então substituir as semanas relativas por datas.}
